% !PRINTED 396
\weberpage{396}
%<<CONT
$Q_n : c$, and to determine $t$ from (5), or also, after $u$ has
been determined, directly from (7).

If $\nu$ is odd and $n$ is an odd multiple of $\nu$, then
\weq{8}{t^2 - D u^2 = -4,}
but if $\nu$ is even or $n$ is an even multiple of $\nu$, then
\weq{9}{t^2 - D u^2 = 4.}
Equation (9) has the self-evident solution $t = 2$,
$u = 0$; in all other solutions of (8) or (9), $t$ and
$u$ are different from zero, and we consider here only the positive
solutions.

We shall now show that from an arbitrary
equation of the form (3), by the formulae (4), (5), one obtains \gesperrt{all
solutions} of (8) and (9).

So let us assume that $t$, $u$ is any positive solution
of equation (8) or (9), and that $\omega = \{a, b, c\}$ is a reduced irrational number belonging to the discriminant $D$. By the equations (6) we determine $P_n$, $Q_n$, $P_{n-1}$, $Q_{n-1}$, which evidently become
integers, since $t$ and $u b$ are either both even or both
odd. From equation (3), which $\omega$ satisfies, (2) then results,
from which in turn
\weq{10}{\omega = \dfrac{P_n \omega + P_{n-1}}{Q_n \omega + Q_{n-1}}}
is to be inferred; and from (8) or (9) it results that
\weq{11}{P_n Q_{n-1} - Q_n P_{n-1} = \mp 1,}
in which, as also in the following formulae, the upper or the
lower sign holds according as $t$, $u$ satisfies equation (8) or (9).
But now from (6) it follows that
\[
Q_n - Q_{n-1} = \frac{(2 c + b) u - t}{2},
\]
and because $\{a, b, c\}$ is a reduced number, we have $2 c + b > \sqrt{D}$
[\S.\,125, (4)], hence
\[
Q_n - Q_{n-1} > \frac{u \sqrt{D} - t}{2} = \frac{u^2 D - t^2}{2(u \sqrt{D} + t)} = \frac{\pm 2}{u \sqrt{D} + t};
\]
but since $u$ and $t$ are positive integers and $D > 1$, we have
$u \sqrt{D} + t > 2$, hence $Q_n - Q_{n-1} > 0$ for the upper and $> -1$
for the lower sign.

Now because $Q_n - Q_{n-1}$ must be an integer, it is
accordingly equal to or greater than zero, and we have
%CONT>>
